% =====================================================================
%  R. Nielsen, TWO KINDS OF INSTRUCTION, I-II
%  In: For Idea and Reality.  A Journal, edited by R. Nielsen, Bj. Bjørnson and Rud. Schmidt.  Volume for 1869.
%  Copenhagen: Chr. Steen & Søn; printed by Louis Klein.  pp. 297-323, pp. 409-441.  (To Slags Undervisning, in For Ide og Virkelighed, Aargang 1869.)
%  ENGLISH TRANSLATION, from transcription.tex (Danish).
%  GENERATED by ebook-method/for-idee-1869/tr_build.py --write from this template and the translators' parts in
%  ebook-method/for-idee-1869/tr/ (TRANSLATOR-BRIEF.md; terms: GLOSSARY.md, the titles and four key terms approved by
%  Hans 2026-10-13).  Conventions: TRANSLATION-PLAYBOOK.md.  Page numbers of the 1869 print in the margin.
% =====================================================================
\documentclass[12pt,a4paper,openany]{book}
\usepackage[utf8]{inputenc}
\usepackage[T1]{fontenc}
\usepackage{libertinus}
\usepackage{libertinust1math}
\usepackage[english]{babel}
\usepackage[protrusion=true,expansion=true]{microtype}
\usepackage{setspace}
\usepackage[margin=1.3in]{geometry}
\usepackage{fancyhdr}
\usepackage{hyperref}
\hypersetup{pdftitle={Two Kinds of Instruction}, pdfauthor={Rasmus Nielsen}, hidelinks}
\setlength{\headheight}{15pt}
\pagestyle{fancy}
\fancyhf{}
\fancyhead[LE,RO]{\thepage}
\fancyhead[C]{\textit{R. Nielsen: Two Kinds of Instruction (1869)}}
\setstretch{1.3}
\usepackage{marginnote}
\renewcommand{\thefootnote}{\arabic{footnote})}   % the print numbers each page's note ¹) (p. 553: *))
\newcommand{\opage}[1]{\marginnote{\footnotesize\textit{[#1]}}}
\newcommand{\essayhead}[2]{\begin{center}{\large #1}\ifx&#2&\else\\[0.3em]{\normalsize #2}\fi\\[0.5em]\rule{3em}{0.4pt}\end{center}\medskip\noindent}
\newcommand{\signature}[1]{\par\smallskip\hfill #1\hspace*{2em}\par}

\begin{document}

% the volume's title page (PDF 7)
\begin{titlepage}
\begin{center}
\vspace*{3em}
{\Huge For Idea and Reality.}\\[1em]
\rule{8em}{0.4pt}\\[1em]
{\large A Journal,}\\[0.8em]
{\small edited}\\[0.4em]
{\footnotesize by}\\[0.6em]
{\large R. Nielsen,}\\[0.4em]
{\large Bj. Bjørnson and Rud. Schmidt.}\\[1em]
\rule{8em}{0.4pt}\\[1em]
{\large Volume for 1869.}\\[1em]
\rule{8em}{0.4pt}\\[4em]
Copenhagen.\\[0.4em]
Chr. Steen \& Søn, Publishers.\\[1.5em]
{\small Printed by Louis Klein.}\\[3em]
{\small Translated from the Danish}
\end{center}
\end{titlepage}

\essayhead{Two Kinds of Instruction.}{I.}
% --- p. 297 ---
\setcounter{footnote}{0}%
\opage{297}Instruction and cultivation can be planned and conducted in very different directions; in attempting
to enumerate all possible methods one would arrive at a number beyond survey. Here attention
shall be drawn only to two schools unlike
in design, aim and procedure,
the \emph{learned} school and the \emph{popular.}

The hallmark of learned instruction is the imparting
of thoroughly coherent knowledge
with a corresponding cultivation of the mind; memory and understanding are
the faculties of the soul on whose development it chiefly
depends: here the work is directed originally at the intellectual
nature of man, not at the whole man. The task
of popular instruction, by contrast, is a personal
cultivation, an influence on the mind, an awakening of the powers of the soul
with a corresponding enlightenment about the goal and conditions of human life
for the individual, for society, for
the race: here the work is directed radically at the whole man.
In learned instruction the interest in knowledge is and must
be the highest, whereas the original driving force of popular enlightenment
is to be sought in the interest in life. As the
interests are, so must the endeavors be. It is
% --- p. 298 ---
\setcounter{footnote}{0}%
\opage{298}necessary to separate the two kinds of instruction properly
from one another when one wishes to gain clarity about the principles;
but this is not to deny that between
the two there must be manifold links, that instruction in knowledge,
for example, can exert an influence, though a mediate one,
yet a considerable one, on character and
direction of life, just as popular enlightenment, for its part,
must, in order to reach any height, both presuppose
and promote the acquisition of the necessary preliminary knowledge.

Learned instruction is among all civilized peoples
in full swing and full recognition. The school of knowledge
is old; through centuries it has known how to
advance with the times, to reform itself, and in spite of
manifold hindrances both internal and external to fight its way
to an ever wider circle of diffusion. ``Knowledge
is power,'' this truth stands clear before all
cultivated people, equally clear before the industrious citizen
who looks practically to the application and
utility of knowledge, and before the theorizing researcher, whose ideal
goal is knowledge itself. Such a recognition
is as yet far from having fallen to the share of popular instruction.
Nor is this instruction centuries
old; on the contrary, it is very young, still in its first,
groping beginning; the vital thought in its inner peculiarity
has --- if only we could say it without
all too painful offense to those who have the disposition for it
--- first of all broken forth with the Grundtvigian
high schools.

Before we nevertheless attempt a closer clarification of
what and how much lies in this assertion, we must
first ward off some misunderstandings which in recent
% --- p. 299 ---
\setcounter{footnote}{0}%
\opage{299}times, perhaps particularly on the occasion of the awakening popular
striving, have arisen concerning the learned system of instruction.

From a somewhat peasant-like point of view people have been
inclined to regard the University and the learned schools
with everything that belongs under them as antiquated arrangements,
institutions whose imminent dissolution must
be considered a necessary condition for freer progress in
the spirit; in complaints about the ``aristocracy of books,'' ``the
learned guild system,'' ``the bread-and-butter system'' and so on, it has been
maintained that the reading of lessons, exercises in composition, and strict examination
transformed the school into a kind of house of correction, an official
institution of compulsion, which cowed the courage of the young, dulled the
innate desire to learn, weakened the vital powers, brought ice-cold
into the warm hearts and brought upon them a spiritual
chill. Although many of these objections,
when we judge by outward reality, are not exactly
wholly unfounded, they are yet, when we look to
the idea, groundless.

The school of knowledge must, in order to solve its tasks,
necessarily proceed from definite presuppositions and
make particular demands. The first condition is that
the pupil impress upon himself what is explained and prescribed to
him by learning certain lessons; no other way to the
acquisition of knowledge is given. Now such
a reading of lessons, especially when the lessons gradually
become more numerous and larger, may often be felt as an oppressive
burden; here there can be talk of lightening the
burden for the young, if it becomes too heavy for him, but here
there can reasonably be no talk of relieving him of
the burden. Burden indicates work, and the young person must
% --- p. 300 ---
\setcounter{footnote}{0}%
\opage{300}learn to work. As now the school on the whole will watch
that the pupil is, without overexertion, put in a position
to work with seriousness, diligence and zest, so it will
also, by careful hearing of lessons, guidance and correction,
know how to control the work. The pupil arrives
in this way at a clear conception of what he
can do and what he cannot do; he learns what it is to
learn something; he gradually perceives more distinctly how
necessary it is to know something thoroughly,
to know it in coherence, solidly, surely and completely.
Thus he comes more and more under the discipline of thought
and feels the weight of valid grounds --- an indispensable
element in the development of character; for the human being who
has no respect for clear thoughts and valid grounds
lacks something essential in being a human being, however living
a feeling, however warm a heart he may otherwise have.

As a complaint against the learned system of schooling in its
present form one not seldom hears the objection that
the program has too many headings, that far too
many heterogeneous things are demanded, so that the pupils learn a little
of everything without learning anything thoroughly. We shall not here
inquire whether perhaps at present too
many subjects are demanded, whether these subjects do not perhaps stand
merely externally coordinate, instead of being subordinated under
a few principal subjects which could give the whole greater consistency
and closure; we shall only point out
that the school of knowledge must secure its intellectual foundation,
and that this foundation becomes too narrow and the horizon
too limited when the subjects of study are too
few and uniform. The scientific preparatory school
cannot begin as a vocational school; for that very reason
% --- p. 301 ---
\setcounter{footnote}{0}%
\opage{301}it takes up so many heterogeneous subjects. The
many subjects mean that it does not depend especially on
any single one by itself, that the single subject is only a single element
in the whole cultivation. Now it is of course quite natural
that among several awakened and gifted pupils one will
have an eye and a predilection for one subject, another for
another, one for history, another for mathematics, a
third for languages, and so forth; such expressions of individual
predilection can even be counted among the
very promising signs. But even if the school can to a certain degree
overlook the intellectual inclinations of individuals,
it dare not, out of regard for the even and, as far as possible, harmonious
development of the faculties, give youthful
amateurishness untimely nourishment, any more than it dares
accept the excuse that a particular head should
be so exclusively organized for a single branch of instruction that
on this account it could not possibly enter
into the rudiments of the others. When the demand is
``to learn something thoroughly in the many subjects,'' this
can reasonably be understood only to mean that the pupil little by
little comes to a cognition of how important it
is to him to work together what is heterogeneous, so that
the overview of the whole stands in the right relation
to the insight into its particulars; it is ideal totalities,
a realm of laws, a world of cognitions, which
gradually is to dawn upon his inner sight. The much-disputed
system of examination, when seen in
this light, will then precisely contribute its share to the attainment of the end.

It is a great and significant task that the student
must solve when he is to prepare himself thoroughly
for a learned examination. It is not done with his
% --- p. 302 ---
\setcounter{footnote}{0}%
\opage{302}stuffing a multitude of pieces of knowledge and doctrines into
his memory; he must understand how to work through the rich,
highly heterogeneous material until he has the whole mass
organized. This requires more than mechanics of the understanding,
memory and diligence; this requires first and last
a wakeful, active, independently working thinking.
True, thinking cannot yet show its independence
by the student's discovering on his own
either new facts or new truths; he is
directly confined to receiving and appropriating what is handed down.
But if we look more closely and consider
what lies in this, we easily discover that the powers which
essentially condition scientific productivity are already
here in full activity. It is, as is well known, by no means
an easy matter to enter completely into
a given doctrine, even when the doctrines are presented
in clear systematic coherence. The difficulty increases
when several mutually conflicting views are to be
judged and what conflicts is to be reconciled; it takes suppleness of spirit,
self-possession and resignation to take possession of
foreign thoughts, to live oneself into them and
make them one's own. The practice of this is so far from
weakening scientific self-thinking that it on the contrary
supplies one of its first indispensable conditions.
How much immature criticism, how many loose, impudent, saucy
dismissals of more profound works are not
brought to market in the literary marketplace solely for the reason
that the critics have neglected that beneficial practice and
for that reason lack the self-criticism which should warn them
against taking up tasks which they have so little worked through
that they would, even as far as the most meagre
% --- p. 303 ---
\setcounter{footnote}{0}%
\opage{303}account is concerned, come off miserably, if
they had to pass a public, impartial, but strict
examination.

The next thing the preparation demands is to get the
heterogeneous cognitions grouped in consciousness in such a way
that the details stand out distinctly without harming
the overview. In this endeavor each person is in the
strictest sense referred to himself; for the procedure
that suits the one may be highly
inconvenient for the other. The psychological interplay
between imagination, understanding and memory, on which
it here depends, can, in its inner, fine, often imperceptible
transitions, be highly different in different people.
He who, for example, has a very comprehensive and faithful memory
will be inclined to begin with the particulars and
regard the whole as a sum of manifold parts; only
by remembering is he helped to think: the understanding operates
on the foundation of memory. He in whom understanding
and thinking have a natural preponderance must go the
opposite way. For him it is first of all important to
bring out the general, to penetrate into the coherence of the propositions
and to preserve the overview. In order not to be burdened
with superfluous material, he must at every step
forward let thinking strictly test and exactly weigh
which and how many particulars in the given
context necessarily belong to the understanding of
the whole: it is the understanding which by this economy
comes to the aid of memory. That both, he
with the strong memory and he with the penetrating
understanding, must strive toward one and the same goal,
goes without saying; but it makes, as
% --- p. 304 ---
\setcounter{footnote}{0}%
\opage{304}regards the manner of working, an essential difference whether it is memory
that originally carries the understanding, or the reverse.

A scientific passage through a great edifice of doctrine, rounded off in
itself, whose portals and pillars, whose architectonic
members and parts stand in symmetrical harmony with
one another and with the whole, and along with it a corresponding
psychological insight into the intellectual nature of the self,
a thorough acquaintance with one's own spiritual-individual
nature, a self-criticism that reveals deficiencies and
excellences and furnishes a measure of future capacity for work,
are advantages which hardly anyone who by a thorough
preparation for examination has endured the strain
will regret having acquired.

Hence, so long as we esteem science and hold fast to
the scientific standpoint, no well-founded
objection whatever can be made against the school as a school of knowledge;
it is much rather to be wished and hoped that
learned instruction, unassailed by the changing
moods of the time, may in one and all remain true to its character.
In the prevailing notions of knowledge
and cultivation, on the other hand, there is a fundamental error which must be
combated with the greatest vigor.

It is a prejudice deeply ingrained in many that the
multiplicity of knowledge is the hallmark of cultivation,
that he who has many pieces of knowledge, can speak many
languages and join in the talk about everything, is the truly cultivated, whereas
he who has less knowledge and speaks only one
language betrays a lack of cultivation and as a limited
individual may easily be overlooked in cultivated society.
That this way of looking at things, which stands in close connection
with the prevailing inclination of the age to
% --- p. 305 ---
\setcounter{footnote}{0}%
\opage{305}place the exclusive superiority of man in the intellectual,
is to the irreparable harm of the life of the spirit in all directions,
can easily be seen upon a little closer consideration.

In the first place it is harmful to learned instruction;
for it brings it about that the school of knowledge incessantly
runs the danger of falling out of its original character.
If it is learning and science that are essentially to condition
the development of the spirit, then the school must most strictly
inculcate that these high intellectual goods are first of all
pursued for their own sake without any
regard to external, finite advantages. To let
oneself be admitted into the learned school is then the same as to
dedicate one's life, one's best time and one's best powers to
the cultivation of science. To be sure, the state ought to see to it
that science can always secure
to its true devotees a real station in life, an honorable
place in society; but such promising prospects must not
come into consideration when the question is of the first
motives. Before science can serve as a means,
it must first be an end in itself, be attributed value
in and for itself and be loved for its own sake. The
disciple who, as he enters the school, does not bring with him
a youthfully great, if also youthfully unclear,
conception of the glorious, the sublime in striving
after knowledge and insight, him the Muses will not greet:
he enters as one unworthy; and the disciple who,
although the school applies its best teaching powers to
open his eye and awaken his sense, yet cannot
get further than to the practical question of what
use there is in learning so many useless things,
ought, even if he otherwise reads the lessons ever so
% --- p. 306 ---
\setcounter{footnote}{0}%
\opage{306}irreproachably, to be given consilium abeundi, for he is
not in his right place, he has gone astray. But if
the school were to maintain the ideal demands in this way,
then it follows of itself that, in order to
gain in content, it would have to lose considerably in extent. To
stay with classical cultivation, two or three
Latin schools, each annually sending on average ten to twelve
students to the University, would probably be about the maximum
in so small a country as Denmark. Now, on the contrary, as unclear
notions about the obstacles and conditions of the life of the spirit
in part decide the issue, the matter stands otherwise.

In the decent opinion that the studious path is after all
always a safe and honorable path, that a treasure of knowledge
is a precious treasure, and that ``the best one can
do for one's children is to give them a good upbringing,''
parents and guardians see themselves concerned about higher
institutions of cultivation, and often see wrongly; for there is
a rush upon the Latin schools, the gymnasia and the University.
This points, especially when we look toward a
coming time in which hardly any well-to-do citizen
will have cause to envy the official class the
golden horn of plenty, evidently to a more deeply lying
fundamental error. It cannot possibly be regard
for the advantage or the good prospects that here has the preponderance;
it must be something else. It is the foggy
mixture of two kinds of cultivation, the intellectual and the
personal, that confuses the conception. --- The Realskole, for example,
is unequivocally a school of knowledge and nothing else. He
who wishes to put his son in a Realskole easily becomes aware
of what it is he really wants. He wants the
young person to work his way forward, to acquire
% --- p. 307 ---
\setcounter{footnote}{0}%
\opage{307}a sum of useful knowledge in definite subjects. He
may otherwise have the greatest respect for scholarliness,
particularly for the real sciences; he may understand
that working with the head strengthens the faculties and gives
the able man a peculiar superiority of understanding; but he
is not easily befooled by unclear fancies concerning the
ideal in the instruction itself: that the Realskole should in
a special degree ennoble the heart and elevate the spiritual
in the development of personality, hardly
occurs to him. Otherwise with classical instruction; it
seems once and for all to have received a privilege to foster
the noblest and deepest in the human spirit,
so that it both imparts useful knowledge and at the same time
gives the young souls the finer spiritual nourishment
which is the condition of ideal personality. But precisely
this divided, two-sided regard, the regard partly to the
useful knowledge, partly to the finer spiritual nourishment,
brings with it an ambiguity that is extremely dangerous to
the school. It dare no longer regard its pupils
exclusively as a circle of young future devotees of science;
far from it! here many are present who notoriously
have never thought of sacrificing their lives to the sciences;
they are here merely for the sake of collecting knowledge
and to obtain some classical cultivation. Still less
dare the school lay the decisive stress on the utility of
knowledge; for as regards the useful, it cannot
possibly compete with the Realskole. If Latin and Greek
are to be learned for the sake of utility, if Herodotus and Homer, Cicero,
Virgil, Ovid and Horace are to be read in the original languages
solely for the sake of the knowledge and for the sake of practical
application, then it looks sad for the finer
% --- p. 308 ---
\setcounter{footnote}{0}%
\opage{308}nourishment, the classical cultivation and the expected nobility of soul.
Thus the school stands in all ambiguity as a
higher institution of cultivation --- with many grade books
--- without character. With this ambiguity it
cannot, namely, be avoided that the pupils who stream in
come to constitute a mixed population. Around
the truly Called, Minerva's legitimate sons, there gathers
by degrees a larger flock of illegitimate sons, of such
as, since they once and for all are to go ``the studious
path,'' at a distance indeed have a certain respect for the
lofty goddess, and probably also in time will do her
the honor of placing her owl on the cap; but
who for the rest, as far as the scientific is concerned, to put it plainly,
do not know what to do with themselves, neither what the whole
is to mean, nor what in a spiritual sense they actually
have to do. The mood of the majority will then exercise
an influence on the spirit in the school which is not for the
better. A capable administration together with zealous,
able teachers can well give assurance that morality,
order and tolerable diligence are maintained on the whole;
but the wakefulness which conditions deeper understanding,
the ideal sense, the spiritual hunger and thirst cannot
be breathed in. As it is repugnant to have to
eat and eat and taste many dishes when one has no
appetite, so it must also come hard to the young person
to have to read and read and learn the hundred
different things of which perhaps hardly more than one
or two could in reality and naturally, --- ``if
he only had time'' --- interest him.

The more the inward considerations decrease, the more the
outward increase. It is not the spirit in the ancients, not the
% --- p. 309 ---
\setcounter{footnote}{0}%
\opage{309}great thoughts of the world, not the rhythms of the language, not the
clear classical form, that captivates the pupil and urges
him to learn his assignment; no, it is other, more outward
and real considerations, considerations of ``mg and ug'' and of
``promotion at the next examination.'' Spiritual work in
this way acquires a likeness to factory work; one attends
to the most necessary part of the daily task mechanically
and with drowsy diligence, but works at full power in the busy
time; one is constantly oppressed by only one worry, a
quiet worry, which grows as the hour of decision
draws near, the worry of ``passing the examination.''
Under the pressure of school compulsion and the in many ways
burdensome reading of lessons there is nevertheless a Something which
moves in the youthful imagination and shines with
the splendor of ideality, it is the notion of the free
life at the University. ``In school we read only for the sake of the
examination; only when we come to the University do
we begin to study.'' This notion will disappoint
many; it will as a rule disappoint all those who have not
from the beginning felt what it is to
set the inward considerations higher than the outward.
For the first ``practical'' demand that the University makes
of the young students is that they as soon as possible
prepare not merely for one, but for three examinations, of
which the two first would by preference have to be concluded in the course of
the first year, so that with undivided powers one can then
work on the great official examination; ``only when the official examination
is well over can one properly begin
to study.'' If already in school time the regard to the examination
has become an outward regard, and the notion of what
it is to read for an examination, break through and
% --- p. 310 ---
\setcounter{footnote}{0}%
\opage{310}consist in an examination, duly trivialized: what wonder, then, if the old sober, trivial notion
pops up anew, so that the factory system, with its great lessons to be learned by heart,
mechanical coaching, hasty repetitions and
hurried cramming, before one is rightly aware of it, is
in full swing.

It is the innate one-sidedness of the school of knowledge that,
when it conscientiously preserves its original
character and does all it can in the service of the spirit,
it is nevertheless not able to develop the whole human
spirit; if, then, the appearance arises notwithstanding that
the learned high school is the only institution of spiritual cultivation
that can be given, it must sink from its original
level down to serve the merely external, to be an
asylum for empty show, a dead mechanical system of learning, a factory system driven by finite interests, from
whose nausea the Muses flee.

If the unclear notion of ``cultivation,'' ``cultivation of the mind,''
``higher cultivation,'' is already harmful to the
scientific system of instruction, then the harm is far
greater still when we descend and consider the instruction
in those circles of life where there can be
no question of science. All our public schools, the burgher schools
and the common schools, are, without exception,
schools of knowledge, all set up for the learning of lessons, learning by heart
and learning for examinations, all of them arranged
after the learned model. If now these schools were all
conceived from the standpoint of popular real schools,
nothing could be objected to that;
for useful knowledge is necessary for all.

But that is not how matters stand. In the burgher and
% --- p. 311 ---
\setcounter{footnote}{0}%
\opage{311}common schools the express aim is to awaken
the higher life in the young generation, to have an educative effect.
To education there belongs, then, not merely that the children learn to
read, write, calculate, etc.; there belongs to it chiefly that
they are instructed about what is highest and most important for life, that
they are instructed in religion. If we now consider the instruction in religion
in these schools a little more closely, we meet
the ambiguity touched on above once more, but in an
even more deplorable shape. To impart knowledge of religion
and to awaken the religious life, these two dissimilar
tasks call forth two endeavors, of which the
one unceasingly works against the other. ``The hour of religion
is a kind of edifying hour; the religious truths
cannot be imparted either as mere doctrines of the understanding
or as external historical facts. For that reason
Jewish children cannot be instructed together with Christian ones;
they can have all the other hours in common, only not the hour
of religion. The doctrines of religion must come
from the heart and go to the heart; they must therefore be imparted
by a believing teacher, who in the best manner seeks
to awaken and strengthen faith in the young.'' This is the endeavor
in the one direction. But with this is now joined
another endeavor, and it is for the school as
school still more important, it is the endeavor to impart
the knowledge in a certain prescribed extent.

What is religion? ``It is called religion to
have knowledge of God and his will, and of the
manner in which he wishes to be honored and worshipped by us''; these
opening words of the older prescribed textbook are to a high
degree characteristic of the standpoint from which instruction in religion
is conceived in the school. In religion, too,
% --- p. 312 ---
\setcounter{footnote}{0}%
\opage{312}there are lessons, in the hour of religion, too, the lesson is to be heard,
in religion, too, an examination is to be held.
What now is it that can be tested by an examination? Not
the life of faith itself, or how far the young each for himself really
have faith. Such an examination in faith the Lord of existence
has reserved exclusively to himself to conduct.
The only thing that can be tested is the knowledge.
But what is exclusively the object of testing in the examination
is also the highest, most important aim of school instruction, and the endeavor for knowledge thereby gains
for this reason a natural and necessary preponderance over
the awakening endeavor, the endeavor for life.

The foundation of the public, official instruction in religion
is an approved textbook, a short extract from the
ecclesiastically systematic dogmatics: here we have the learned
model. The textbook with its scriptural passages is intended
for learning by heart; it is especially a matter of having the scriptural passages
at one's fingertips. Learning by heart is for children
a spiritually mechanical occupation; so it begins with
spiritual mechanics; but spiritual mechanics in religion strikes
life dead. It is of no use that the textbook expressly
says that religion is not only a knowledge, but
also a life; for, when it stands in the lesson, it passes
instantly into the mechanics of recitation that religion is
a life; it is of no use that a new textbook is authorized
which begins a little more heartily with the question:
``What should above all lie upon our heart?'' for,
what ``should above all lie upon our heart'' runs
along without more ado in the rigmarole. If a
prize were offered for the answer to the question: ``How
shall one, under the appearance of wanting to impress upon children
% --- p. 313 ---
\setcounter{footnote}{0}%
\opage{313}the chief doctrines of revealed religion, contribute without offense
one's part to their being able from the very beginning
to conceive a secret loathing for all religion?'' then the task would
without doubt be best solved by a thorough psychological
development of what, during the learning by heart of a
dogmatic textbook, must go on, half consciously, half unconsciously,
in the child's soul.

It is certainly possible with some justification to appeal
to the fact that it is not a slavish learning by heart that is at issue,
that the textbook is only a guideline for sensible
use, and that the teacher will always take care
to awaken the children's reflection, by explaining to them
what has been read in such a way that it both acts on the feelings
and enlightens the understanding. We do not doubt that the
upright teacher will do his best in both respects;
we only mean that what is in itself contradictory
cannot be carried out. As regards the feelings, it is
impossible that truths which are imparted in the form of mere
doctrinal propositions and have so little to do with the child's real
life should precisely leave any particularly deep
impression or fasten themselves deeply in the childlike mind. Let one
take, for example, a doctrinal proposition like this: ``Marriage
is the intimate union between man and wife,
which God has instituted for mutual help and the propagation
of the human race. Married people shall live together
in love and fidelity until death.'' However true this
proposition may be in and for itself, it would yet
be unnatural if it acted strongly on the little
ten-year-old girl or awakened serious thoughts in the eleven-year-old
boy. If the little girl and the little boy,
when outside school hours they talked with one another about what
% --- p. 314 ---
\setcounter{footnote}{0}%
\opage{314}naturally must lie on both their minds, felt
themselves involuntarily prompted to go through and explain
the said passage, together with the scriptural passages belonging under it,
with one another off the cuff: then one would surely
find it somewhat questionable, indeed uncanny, because it was
unnatural. On the other hand one would find it quite in order
if the diligent children, the little boy and the little
girl, in reading hours outside school heard each other recite
the prescribed lesson and kept careful watch over the scriptural passages,
both the one with ``Let marriage be honorable'' and
the one with ``Husbands are indebted'' and the one with ``Ye wives!
be in subjection to your own husbands''; the childlike thing
is that nothing further is thought in it, that it goes
by rote; the faster the rote goes, the more innocent
is the whole.

The textbook is objective; its dogmatic propositions
are all arranged according to the system, and the objective system
cannot take account of the as yet immature
subjectivity of the childlike age. Let no one object that it is only
a few doctrinal propositions whose application and right understanding
lie outside the children's range of view; for it is easily
shown that there is not a single religious proposition
which, mark you, in the form of a textbook proposition, is, if
we take it strictly, accessible to the childlike understanding.
``God is just, that is: he lets the good
prosper and be rewarded, and he checks and punishes the
evil''; this proposition a child cannot understand. If it is
taken simply according to the words, it comes into conflict
with facts of which the young person will become aware even early
in life; if it is taken simply according to the spirit,
then the meaning becomes this, that the truly good,
% --- p. 315 ---
\setcounter{footnote}{0}%
\opage{315}as may be seen from the example of the Prophets, of the Apostles and above
all of Christ, precisely must suffer in this world,
while evil in many ways gets ``prosperity
and reward''; but up to that height of the spirit,
from which so cutting a disproportion shows itself compatible
with God's justice, the child cannot be raised.
If the edifying interpretation is to aim at
showing that evil may indeed for a time prosper
in the world, but only for a time, so that, considering
the outcome, it generally pays best to
do the good: then the explanation becomes reasoning;
there arises a mixed concept, a hodgepodge of
piety and self-interest, fear of God and worldly prudence,
which confuses the child's thoughts and strikes down its ideal
conception of God's justice.

``But'' --- one will perhaps say --- ``this whole way of looking at things
produces its effect only by exaggerating
the shadow sides and staring at the imperfection
that clings to everything human; it overlooks the facts of
life and stands in conflict with sure experience.
For experience teaches that decent, moral, diligent,
reasonably gifted children without difficulty both learn
the textbook and understand the explanation and find edification
in it. Anyone who doubts this need only, when opportunity
offers, step into the school and attend a
catechization or, what is still more reliable, be
present at the actual test of maturity, the final
public examination. By listening to the examination every
fair-minded person will be able to convince himself that the children
really have understood the textbook, that by means of connected
knowledge of the whole of religion a
% --- p. 316 ---
\setcounter{footnote}{0}%
\opage{316}sure ground has been laid, so that faith is certain and the understanding enlightened.
It is surely a most important matter that the young people's understanding
of religion is developed early, so that they can give
an account both of what they are to believe, how much they
are to believe, and why they are to believe it all.''
We understand very well what lies in this. It is with
an often astonishing assurance that awakened and well-instructed
children can at a school examination answer questions
that otherwise put graying thinkers in embarrassment.
``What is meant by saying that God has created heaven and earth?''
``Are there other rational creatures besides men?''
``Are all the angels good?'' ``What is God's
purpose in sustaining the world?'' to such questions
the half-grown child can often answer so clearly, purely and
definitely that it puts its elders to shame. If now
there really were truth in all this understanding, then
the marvelous would be attained. But the truth is that the whole
is show, a psychological, soul-corrupting deception.

Everyone who has taken any trouble with teaching
will have become aware of the suppleness
with which a youthful understanding is able to imitate
the teacher's thoughts without actually thinking them. There is
an imitation of thoughts that can be so deceptive
that it looks to the life as though the pupil had really
understood the matter, in that he even in part in his own
words exactly reproduces a prescribed movement of thought;
and yet it is a deception, a for the most part naive self-deception. In the exact sciences, mathematics, mechanics,
the deception is still fairly easy to detect; but not
so in the doctrine of spirit. In philosophy, for example, one will be able to
present any problem, even the very most difficult,
% --- p. 317 ---
\setcounter{footnote}{0}%
\opage{317}in such a way that the beginner becomes able to
follow the exactly marked path of thought through all
its windings with such clarity, circumspection and apparent
naturalness that it looks as though
the young man had thought about the matter for many years; and
yet the whole is perhaps only a deceptive imitation.
Therefore in philosophical instruction care must be taken above
all that the scheme does not come to count
as the system, that a stereotyped textbook with measured-out
chapters and standing paragraphs does not, year after year,
settle so firmly in tradition that philosophy itself becomes
a lesson in thinking. That a certain thoughtfulness is required
to learn a lesson in thinking is surely certain; but it is
also certain that there is nothing that weakens more the capacity
for thinking for oneself than this hollow, superficial imitation,
this thoughtless thinking, which in the end must issue in intellectual
mendacity.

If thoughtless thinking is a dangerous parasite
in scientific instruction, where so many and different correctives
can after all be applied, then it is,
as everyone will understand, still far more dangerous in
the instruction of children, most dangerous of all in instruction in
religion. The objective religious doctrinal propositions lie to such
a degree outside the children's real power of comprehension that these
can neither feel nor think anything truly essential
in connection with any of them. Take, for example, the doctrine of the Lord's Supper.
``The bread and the wine in the Holy Communion are''
--- says the textbook --- ``not mere signs of the body
and blood of Jesus; nor are the bread and wine transformed
into the body and blood of Jesus, but in, \emph{with} and \emph{under} the bread
and wine we receive in the Holy Communion the
% --- p. 318 ---
\setcounter{footnote}{0}%
\opage{318}body and blood of Jesus, which in a manner incomprehensible to us
is really present there.'' We ask: what impression
is this dogmatic enlightenment to make on the young
girl's childlike mind? How is the half-grown
boy to be able either to think or to imagine
anything definite by the given explanation? The most natural thing
is that he does not even try.
Should it occur to him for a moment to want to understand
the explanation, then he is badly off. ``Whence do we know
anything about the manner in which the body and
blood of Christ are present in the Supper? It must surely be a
divine authority that tells us this; but which?
It cannot be the Bible; for from the words of institution:
\emph{this is my body; this is the New Testament's
cup in my blood,} we cannot see how
things stand with the \emph{manner.} For it does not
say: in, \emph{with and under this} (the bread) \emph{is my body;}
it says expressly that \emph{this} (the bread) \emph{is}
\emph{my body;} still less does it say that the blood is in
the wine, for it lies quite as near to understand the words of institution
thus, that it is the \emph{wine} (the New Testament's
cup) \emph{that is in the blood.} If we do not have a
divine authority, then we must surely have a human one.
It must be great teachers of religion, learned and
pious men, who by earnest thinking have fathomed
the manner, and so definitely and exactly that we can now
as safely rely on their human authority
as on the divine authority of Scripture; but is that
possible? Here it is expressly taught that the manner is incomprehensible;
how then is one to imagine that there are
men who can exactly establish and determine,
% --- p. 319 ---
\setcounter{footnote}{0}%
\opage{319}therefore comprehend an incomprehensible manner. Here one is both
to think and yet not think, both to imagine and yet
not imagine anything definite by \emph{the incomprehensible}
\emph{manner:} how strange!''

It is a great good fortune for the teachers of religion in the school
that it does not even occur to the best-gifted boys
to want to think over what stands in the textbook.
If this bad habit got started, it would cause much
disturbance in the lessons; for even if the children were
so morally sincere that they never allowed themselves to ask
either from self-importance or from random whims: still
there would come so many questions that---for the sake of
the recitation---there would not be time to answer
half of them, and, what is worse, questions that
would put the teacher himself in the most painful embarrassment.
As it is, however, since that bad habit fortunately does not occur,
the matter goes its even, quiet course. The young learn by
instinct, for instinct is a fine and acute teacher,
that in the impressed notions there is something one is not
to imagine, in the understanding something one is not to understand:
so the young learn, not to think religiously, no, to imitate
and reproduce prescribed religious lessons in thinking, so that
the talent for imitation with its brilliant hollowness can come to
honor and esteem at the examination. That is not true cultivation,
not even the beginning of it; it is mechanical
drill. As the recruit is drilled to learn his regulations,
so the growing schoolboy is drilled
to learn his dogmas; and yet not even so,
not nearly so practical, for the young soldier will soon
have use for the regulations, but the boy is still
far from having use for the dogmas.

% --- p. 320 ---
\setcounter{footnote}{0}%
\opage{320}If the interest of life were really set above the interest of knowledge,
then it would not be the textbook but biblical history
that came to furnish the foundation. Now biblical history is
indeed taught in the common schools and the burgher schools alike;
but in this respect, too, the learned model prevails.
Instead of letting the great events of the spirit
come forward to intuition through the Bible's own childlike
exalted presentation, in which the opposition of the divine
and the human, of the heavenly and the
earthly, the ideal and the trivial reality, can
be felt and grasped without reasoning circumlocution, there is used
now, for a schematic survey, a little spiritless extract of
the whole history of religion, so that the child may equally get
to know something about all the things God has done from the beginning of the
world, from the first day of creation down to the Diet
at Copenhagen in 1536. Here there can be no dwelling on
single great scenes, on such thought-laden and colorful
pictures of life as might give the young souls a deep, unforgettable
impression of the glory of the ideal, of the power
of spirit, when the voice of God is heard, when the marvelous
happens, when the supernatural shines upon the natural.
For the sake of the extract all events, the supernatural
along with the natural, are set on the same level of reality
and placed in the same light: that ``the Lord sent his
angel to Daniel's three friends in the burning furnace'' is
real, that ``Jonah was swallowed by a great fish, which
cast him up again, so that he came alive onto the
shore'' is real, real in the same sense
as it is real that Rehoboam was contemporary with
Jeroboam, that Titus was bishop in Crete and Kristen
Pedersen was born in Svendborg. So one has an even
% --- p. 321 ---
\setcounter{footnote}{0}%
\opage{321}overview of everything, gets to know in a certain connection a little about
every kind of reality, a little about Titus, a little
about Daniel, a little about the burning furnace, a little about the
great fish, a little about Jonah and a little about Kristen Pedersen.

It is of course no use to want to reason with
children about the difference between the true and the real,
between reality of a lower and of a higher order;
but a little more caution in stringing together the many
miracle stories and presenting the supernatural
in so sensuous an external manner that the first
step on the path of experience and of thinking for oneself must
lead the young into a dangerous temptation, might surely
be exercised.

It is not proposals for improvements in the existing
school system that we put forward here, any more than
it is the introduction either of new textbooks or new
biblical histories that we recommend. On the contrary, we proceed from
the presupposition that the school, as it now is, deserves all
recognition---as a school of knowledge. Only this we must
emphasize in the most emphatic way, that so long as there are,
even as regards the instruction in religion, no
other popular schools than the schools of knowledge arranged
after the learned model: these schools of knowledge,
be they otherwise as excellent as they like,
can never lay the foundation either of true personal cultivation
or of ``living popular enlightenment.'' In every fruitful
instruction either the interest of knowledge or the interest of life
must be the decisive one; to want with a certain evenness
to satisfy both requirements and, as one
says in plain speech, to kill two birds with one stone,
is to want the unreasonable, the self-contradictory, the impossible.
% --- p. 322 ---
\setcounter{footnote}{0}%
\opage{322}So long as the acquisition of knowledge is the goal of cultivation,
and the multiplicity of the knowledge acquired the measure of cultivation,
so long will the great majority of the population
in every country live on in an unfree, commonplace condition;
for as regards the acquisition of knowledge the
majority cannot possibly compete with a minority of the more
favored. But if the great majority by lack of
knowledge is not in a position to attain a true, personal,
independent cultivation of the mind, then the future, in spite of
political rights and free constitutions, will not become
a time of rights, equality and freedom, but on
the one side a time of the refinement of culture, on the other
of coarseness, and in both respects a time of lawlessness,
confusion and bondage; for external
freedom can be grounded only on the internal. ``There is
and can be a free people, but there is not and cannot
be a free common populace; for the stage of the common populace is, spiritually
seen, that of minority, hence of unfreedom. What matters
in popular instruction is to bring the common
man so far that he will be able to think clearly his own
thoughts, to give himself sufficient account of his own actions;
to see and understand that no one in society can
have rights without also having duties; to see
and understand that the interest of the individual is promoted only in a
rational way by the interests of society being looked after
on the large scale; to see and understand that the interests of society on the large scale cannot possibly be promoted unless
the various classes of labor, not only those who
work in the bodily, but also those who work in
the spiritual, enjoy full recognition, each according to the conditions
grounded in the nature of the circumstances. Only he who
% --- p. 323 ---
\setcounter{footnote}{0}%
\opage{323}is in a position to give himself grounds for what he wills
and what he does will be in a position to understand the
grounds that are given him by others. The difficulty in
the emancipation of the common man is not that he lacks
either noble self-respect or understanding or honest will;
but the difficulty is that his thinking is unclear, his
range of view limited, his mind suspicious. The difficulty
is that he does not rightly know either how to give grounds
or to receive grounds, because he is essentially unfree.
The common man is unfree so long as his interior is obscure;
but that his interior is obscure makes itself known thereby. that   % print: derved. at (as printed)
he is in the most personal concerns dark,
closed. --- There is in the people a double need, need
of religion and need of enlightenment, need of the redemption of the soul
and need of culture''\footnote[1]{On Obstacles to and Conditions for Spiritual Life in the Present Age. P. 287---88.}; if this need is to be relieved,
then a start must be made afresh on a new path,
with a new instruction, for centuries, indeed millennia,
of experience have taught that the goal is not reached with
the old one.
\signature{R. \emph{Nielsen}.}

\clearpage
\essayhead{Two Kinds of Instruction.}{II.}
% --- p. 409 ---
\setcounter{footnote}{0}%
\opage{409}If the difference between the new schools and the old consisted
solely in this, that the former offered directly a continuation
and extension of the instruction that had already been
founded in the latter, then the hoped-for progress would
surely not amount to much. Were one simply to
build on the old foundation, one would also have to employ,
in all essentials, the old method. Even if a new subject matter
were chosen, in that the peasant high schools aimed in particular
at making the common man acquainted with the legendary history,
literature and political conditions of the fatherland, that would
still not change the matter. For how was the method to be arranged?
To make the communication of knowledge a little freer,
livelier and more entertaining through the exclusive use of the
spoken word and by exempting the pupils from the reading of lessons
and from examination, would hardly lead to the goal; the consequence of such
a freehand experiment would probably rather be that the
new disturbed the old without itself coming into strength.
No, the truth is that from the school of knowledge to
the popular school there is no directly paved transition.

Rooted in the old, the doctrinaire man of instruction
is placed in the painful embarrassment that he
% --- p. 410 ---
\setcounter{footnote}{0}%
\opage{410}cannot but feel the vital power that breathes in the new, and
yet must stare at this new with a vexed and indignant
look, a look that stops at the shell without
discovering the kernel. From this quarter various objections
will then also naturally be heard. ``One now begins
by making popular instruction a kind of agreeable
pastime, a sort of proverb-game with song and
sound and all manner of edifying amusement. But to work
spiritually is not the same as to amuse oneself. It is
`earnest and play, that are to go so well together'; and then one chooses the play that is fitting,
not for the teacher, but for the nurse. That is indeed to degrade
the divine reason. All true instruction
must begin with reverence for that which is
higher than man, so that everyone who enters
a great institution of learning. can be reminded of what the Word says: % print: Læreanstalt. kan (as printed)
take off thy shoes, for this place is holy. There is a
nursery-like naturalness with which one, freely, by whim
and caprice, dares to fall into chatter about the very highest
concerns of life, and to peep at `Our Lord' in his cards in a chummy way:
is that supposed to be the new? In the self-complacent
conceit that sturdy tastelessness and logical unceremoniousness
furnish the most excellent proof of popular eloquence,
one fills the heads of the young apprentices with ballads
and proverbs, Emter and fairy tales, folk-historical and world-historical % print: Emter (as printed)
surveys, until at last it must make their eyes
swim. Following a conception of spiritual phenomena in
figurative language that is more dreamy than discerning,
one explains image by image, and nods admiringly
at every saying without knowing a thought.
Far from passing for what they originally are,
% --- p. 411 ---
\setcounter{footnote}{0}%
\opage{411}symbols of ancient thoughts, which only in a given context
have a certain restricted universal validity, legends
and myths are treated as oracles of unfathomably
profound content, and their interpretation is made dependent on
the most arbitrary association of ideas. What the ``myth-smith''
meant is a matter of indifference; the important thing is that
the myth can be used for anything whatever: a casket of images,
in which one keeps what the bard's glance has seen, a
coloured wrapper around the prophetic spirit's very latest
historical revelations. It is a false striving
to be at once profound, witty, and yet simple.
To counteract one extreme one throws oneself
into another. One would avoid stiff forms
and becomes formless; one would get away from the artificial
and lapses into flatness; one scorns vain finery,
and takes hold of the pure with unwashed hands. Such
a straightforwardness, whereby one in a chatty fashion familiarizes
oneself with the ideal powers of existence, makes an impression
much as if a congregation in one of the free states,
in order to avoid false solemnity, hit upon the idea of assembling
and holding prayers with hands in pockets and a cigar
in the mouth.'' --- Let this, then, be a sample of doctrinaire
understanding; it is quite natural that those
who are unable to penetrate into the fortress itself must
keep to the outworks.

The fundamental difference between the new and the old is
conspicuous when we look to the first conditions.
The old, official popular school is a school for children. Against
children and minors compulsion may be applied; the children's school
under the guardianship of the state is accordingly also a compulsory school
in a double respect: for the children and for the
% --- p. 412 ---
\setcounter{footnote}{0}%
\opage{412}parents. The new popular school is a school for adults;
for the instruction of adults society has no compulsory laws:
here, then, one can begin with freedom. The old
popular school is a school of knowledge and, as a school of knowledge,
must build on an artificial, mechanical foundation;
the new popular school, on the other hand, which through the awakening of
the personal powers of the spirit will be a school for life,
can speak with living human beings about life itself, about
what lies close to them, and begin with a natural
foundation. By considering these starting-points
each by itself somewhat more closely we shall easily convince ourselves
how essential they are.

It is a circumstance on which weight must be laid
that the official instruction of youth ends precisely
at the point of time when the life of the spirit properly
is only to begin in earnest. That one cannot, by lessons of thought
for rote learning, anticipate the reflection
expected of a maturer age has already been shown. It
is not with wisdom for living as with a confirmation dress, which
can be sewn too long and too wide, so as to fit with the years.
By giving the child a systematic abridgment of
moral doctrine one does not further its morality; still
less does one strengthen its faith in God by a children's dogmatics
composed in textbook form, for children's dogmatists
and dogmatic children will hardly ``meet in the Kingdom of God.''
What holds of religion and moral doctrine holds in
a corresponding way of world history, the history of the fatherland,
the demands of civic life, etc.: there is a limit
to the comprehension of minors that cannot be passed
without a breach of the divine order being made.
In this respect it is of importance that the new
% --- p. 413 ---
\setcounter{footnote}{0}%
\opage{413}popular school begins instruction afresh, in that it begins
with adults.

If it is already in itself unnatural that a school for
adults should condescend to work by compulsion, then
the unnatural becomes still more unnatural when
it is expressly a school of awakening, a school for life.
In the mere acquisition of knowledge, in particular in
everything that comes under the reading of lessons and examination,
it is a rather indifferent matter whether the drilling
is done voluntarily or involuntarily. When a boy has learned
the big multiplication table, then he has learned it, and then the toil is
over. Whether it went with the good or with
the evil, whether it was a true pleasure to him or it
cost tears, he may well forget, so long as he only
does not forget the table. Otherwise with the truths
that are to work upon the mind and have significance for personality.
To be sure, there are in practical life
many actions that are neither wholly voluntary
nor wholly involuntary, but a mixture of both,
such as paying taxes, doing military service, and the like.
Compulsion and freedom can be reconciled in all such cases
where it matters more about the outward side of the action
than the inward, more about the deed than
the motive, on the whole more about society than about
the individuals. But a view of life, personal conviction,
love of the fatherland, living love of God
and neighbor, such gifts of the spirit cannot be collected
by compulsion; they are the fruits of freedom, and freedom
has its beginning in itself.

The concept in which the unity of compulsion and freedom originally
must be sought is the concept: upbringing. Upbringing
% --- p. 414 ---
\setcounter{footnote}{0}%
\opage{414}presupposes discipline and is to that extent conditioned by compulsion;
but true upbringing must always end in self-upbringing.
Discipline by others acquires significance only by
leading to self-discipline; self-discipline demands a firm
will, a vigorous self-determination: here freedom begins.
The two tasks lying in upbringing can now be
equally divided between the opposition of the old and
the new school.

The watchword in the old system of upbringing has always
been, and will, it is hoped, continue to be: respect and
obedience. With this imposing demand the school of knowledge
meets the young, and then makes use of the conditions
and means that instruction places at its disposal,
in order to let them feel that there is something called
law. To be in school is to be under the law.
Strict discipline is not a necessary evil, no,
it is an essential good, for rational strictness is widely
different from despotic harshness. Strictness is the expression
of earnestness, which makes the unconditional palpable and strikes
unruliness with the poet's words:

\begin{quote}
``Respect is what is lacking,\\
It must be impressed upon thy mind,\\
And thou must subordinate thyself\\
To that which is a little more than thou.\\
Respect, as said, for now thou hast none,\\
Respect for the place, respect for the thing.''
\end{quote}

It is presumably above all the respect for the thing, to which the
old school had disciplined its disciples, that the end
of the last century and the beginning of this one
could still thank for so many great public characters of energy.
Character presupposes firm
% --- p. 415 ---
\setcounter{footnote}{0}%
\opage{415}conviction and strong will; but firm conviction and
strong will presuppose in turn that there is something that
unconditionally wills to be respected. Through respect for the thing
the strife between opposing views of life could break through
in the age and become somewhat more than a wrangling over opinions
among newspaper-writing individuals: every view that wished to
stand and fall with its character was esteemed; only
the fickle characterlessness that wills nothing was despised.

Where respect for the thing is lacking, one arranges matters
in another way. On the presupposition that every
definite view, just like every definite character,
is a one-sidedness, a kind of narrowness in the will,
one sets up a literary shop with principles and
theories in all patterns, old and new, positive and negative,
advancing and retreating according to circumstances, a kaleidoscope
of the most contradictory views. Every time
someone shows himself with a new view, there is always a
journeyman in the shop who waylays him as he passes by,
flings in his face, for the amusement of the public, an old and
opposite view, and laughs him to scorn, not
because the passer-by's view is just worse than
so many others', but because he has the presumption,
the narrowness of will, to want to have a definite view
and to resemble something that might be called character.
If in this way it succeeds, little by little, in transforming
the respect of a most honored public for the thing into respect
for the shop, then the game is won. Then an
honorable owner of such an honorable shop may sit snugly
in public and puff himself up in all possible honorableness;
unless there were one honorableness
he could not attain, that he could not attain to become
an
% --- p. 416 ---
\setcounter{footnote}{0}%
\opage{416}honorable character. For character, according to the presupposition,
is stamped as one-sidedness; but stamped one-sidedness
bursts the shop, for the shop is airy
characterless all-sidedness. That an individual who cannot
become a character cannot become an honorable character either
is surely as evident as that a being who
cannot become a human being cannot become a fellow human being either.
``Respect is what is lacking,'' respect,
not for the individual, but for the cause, when the cause is a
question of life, the answering of which cannot possibly be avoided,
unless characterlessness itself is to be the age's
character. Lack of respect for an important cause
is lack of upbringing, a testimony that the
reasoning individuals have withdrawn from the strict
discipline by running away from school too early.

Strict discipline is the honor of the old school;
it stands on the ground of law and points toward the promised
land of freedom, but does not itself enter there; where
compulsion ceases, it has its limit. At this
limit it takes leave of the disciples, who must now
on their own venture the Leucadian leap into the
new region. The transition from the old to the new is
marked by the turning-point at which compulsory upbringing
is to be transformed into self-upbringing. With regard to the
higher popular school this transition is expressly made
dependent on a choice. Not what is chosen, but that
a choice is made, is here the decisive thing. As far as
compulsion extends, the responsibility falls essentially on guardians
and teachers; when compulsion ceases, the responsibility falls
on the person concerned himself. It now becomes his
own affair whether he will be content with what he has, or
% --- p. 417 ---
\setcounter{footnote}{0}%
\opage{417}strive after something better. The school of freedom invites him,
offers him its service, but it does not tempt him.
It has no final advantages to offer him, no
alluring prospects, as far as externals go; it
must, on the contrary, begin by demanding even relatively
considerable sacrifices of money, time and powers of labor.
That so many, even among the impecunious, have already
accepted the terms is a gratifying sign of the
times. It bears witness to what stirs deep in the mind
of the people: a need for awakening, for inner liberation.
The common people sense what lofty cleverness
has not been able to calculate, that there must be a truth
which, independently of finite difference in knowledge and learning,
is accessible to all, a living truth, which
kindles its torch at the sun of freedom and makes spiritual
guardianship superfluous. There is, then, every psychological
certainty that young men who under the stated
conditions enter the school of freedom will not either
misunderstand or misuse freedom. The higher school of knowledge
must always ask those entering about their preliminary knowledge,
for in preliminary knowledge the preparation consists;
the school of freedom, on the other hand, does not ask about preliminary knowledge,
it puts the preparation in the need. Not
the pupil who has the most knowledge, but he who
has the most awakened sense and feels the strongest impulse
to active participation in the spiritual work,
is the best prepared. Nor does he come
for appearance's sake; the sacrifice is too costly for that;
no, he comes for the sake of the cause. He will not drowsily
gnaw on a lesson of thought, unconcerned about the thought
itself, attentive only to how the teacher ``wants
% --- p. 418 ---
\setcounter{footnote}{0}%
\opage{418}it''; he will, as far as it lies in his power, exert
himself to see with his own eyes, hear with his own ears,
think with his own brain: such a fresh originality is
surely able to outweigh the lack of many years of preparatory study.

The foundation for the instruction now presents itself of itself;
for the school begins with the beginner there where he
is, and speaks to him only about things that he can without difficulty
understand at once. To interpret a proverb or
to hear a folk legend, a fairy tale, a piece of history, everyone
is prepared; daily experience has moreover enriched
him with a not small acquaintance with human life
and the nature surrounding him: here are points of connection
enough, if the teacher understands how to use them.
But how are they to be used? Here we meet a difficulty
that we must not overlook, and the less so
as it can easily be magnified by those who wish to
make objections. ``Let the beginning'' --- one says ---
``be ever so natural, and the spoken word ever so
awakening: here is a subject matter, then, after all. The subject matter will grow
every day with every hour: how is the listener
by mere hearing to be able to collect and order the accumulated
subject matter, if he has neither books nor written
notes? But if, in order to profit from
the lecture, he must have recourse to textbooks and written notebooks,
then the method at the folk high school becomes in all
essentials the same as at the learned high school. In order that
the spoken word may bear its fruit, the pupil must
beforehand have learned both to read and to write;
he must begin with knowledge in order to acquire
knowledge: in this way the school for life becomes after all
a school of knowledge; and what difference is there then
% --- p. 419 ---
\setcounter{footnote}{0}%
\opage{419}between the new and the old? That the folk high school has not
been able to maintain itself on the height of sheer awakening
and enthusiasm, but has had to condescend to the
real, and from the beginning to arrange itself to impart
useful knowledge in definite subjects, is a fact.
Thus one finds, e.g., in ``Report on the Folk High School
at Rødding by C. Flor 1846'' timetables
that read ``Theoretical Agriculture,'' ``Physics,'' ``Geography,''
Statistics,'' ``Arithmetic and Reckoning,'' ``Geometry
and Drawing,'' etc. --- Were the strongly emphasized difference
between the old and the new to have validity in principle,
then progress would have to be marked by
the discovery of a new method, a way of instruction
that made it possible to impart to unprepared pupils
thorough knowledge without troubling them with the reading of lessons
and examination. Has such a discovery really been
made? And if it had been made, what then? Then it would
immediately benefit the old school and precisely
serve to level every chasm between the old and
the new.''

That objections of this nature rest on an obvious
misunderstanding is easy to show. In order that the school of knowledge
may answer to its concept, the communication of knowledge must
be end, not means; in the popular school, on the other hand,
the communication of knowledge must be means, not end. That
makes an essential difference. Let the subject be, e.g., geometry.
If one begins thus: Geometry is
the science of spatial magnitudes, a system of connected
theorems; it is therefore necessary to proceed
from proof to proof without skipping any essential
intermediate link; every proof acquires its significance through
% --- p. 420 ---
\setcounter{footnote}{0}%
\opage{420}the presuppositions from which it proceeds, and the corollaries
to which it can lead further, etc., then we hear
at once where we are; we are in the school of knowledge. In the popular school
one begins from another starting-point. It is
not the demonstration in and for itself, not the relations between
mathematical lines, angles and surfaces, that originally
awakens interest; it is as a rule problems of a purely
practical kind. That a triangle is always completely determined
when two angles with an adjacent side are
given will hardly appeal to the beginner to any considerable degree.
``But how is it possible that, standing
on the ground, one can measure the height of a tower or of a
tree in the forest? That problem would well be worth getting
solved.'' By starting from something real, from something that
is immediately familiar to him, the pupil is easily led to an
understanding of what the doctrine of the relation between the parts of the triangle
means. Interest in the demonstration is derived
from interest in the practical problem. The number of
theorems and proofs that he must go through, and the
order in which these follow one another, are accordingly governed
not by the system, but by the problem chosen. It
is not even necessary that all theorems be proved.
When the pupil has once begun to see what
a proof is, many propositions can simply be taken on
authority. At a feast in the country, where the conversation on the
occasion of the almanac had begun to take a
higher flight, someone once asked: ``how far
is it, I wonder, from the earth to the moon?'' ``Nobody
knows that,'' answered an old peasant; ``none
of the learned has yet driven that road.'' --- Were
a pupil without preliminary knowledge to work his way systematically
% --- p. 421 ---
\setcounter{footnote}{0}%
\opage{421}through arithmetic and geometry into astronomy,
until in the chapter on parallax he could get that question
answered, the labor would, in proportion to the short
time he can spend on it, certainly be insurmountable.
And yet the answering of a single such
question can be of extraordinary significance for his
personal cultivation; it can open to him an insight, not
only into the structure of the universe but also into human
knowledge, an insight that raises him high above the plebeian
standpoint and gives him an idea that
science and common enlightenment have something to do with
each other. Fortunately he can, if he is truly
in earnest, with an easily surmountable preparation
get his question answered, without collapsing
under the burden of systematics. The same holds of
the principal proofs on which the Copernican system
rests. Here, by a practical procedure, knowledge can be acquired,
even connected knowledge,
e.g. in land surveying; but the knowledge is means,
not end. The extent and the connection
in which it is imparted depend on the pupil's individual
gifts and personal interest. It is not for the sake of science,
but for his own sake, that he works;
it is not the amount of what he learns, but
the guidance he receives to develop himself, that is
the main thing; it is awakening on which it first and
last depends. ``The most important thing about our high school'' ---
it says in the cited ``Report'' p. 13 --- ``is not the
positive knowledge and skills that we seek to
impart to the pupils, but rather the whole spiritual life
that is awakened and nourished in them, so that their
% --- p. 422 ---
\setcounter{footnote}{0}%
\opage{422}understanding becomes sharper, their power of judgment more mature, their % print: Dömmekraft (as printed)
hearts more open and noble, that there is awakened in them
a sense for order, beauty and a tasteful life, that desire for
occupation may take the place of the usual
habit of sluggishness, that their mind and thereby their nature
may become more cheerful, their feeling for comradeship
and fatherland be awakened, nourished and strengthened.''

If it already holds of the circle of knowledge
that concerns most nearly the external and finite things, that
it furnishes means of awakening for personal self-development,
then this holds in a still higher sense of the elements of cultivation
that by their nature belong exclusively
to the realm of spirit; we think of mythology and history
and the relation of both to revealed religion.

``Whoever wishes eternal war'' --- says N. F. S.
Grundtvig in his presentation of ``The Champion Spirit of the North''
(Mythology of the North 1832, p. 95 ff.) --- ``has an easy time
reaching his insane aim, for he need only
proclaim open feud on honesty and truth, and his
will is done; but whoever wishes eternal peace, and who
should not always wish it? he by no means gets
there so easily as the dear philosophers everywhere have
made them believe, and as one now usually thinks. Peace and
quiet are one thing, and stillness is another, so that tranquillity
can be united with strong movement, and stillness can
be bound up with great unrest, as one can perhaps see most clearly
by comparing the magnanimous, victorious hero with
the slave bound hand and foot. This was also the image that
hovered before our free-born, high-minded ancestors, so
that they felt deeply that life is the first condition of what
a living people can call peace and quiet, and the free
% --- p. 423 ---
\setcounter{footnote}{0}%
\opage{423}development of the powers the other\dots{}. This Northern view of
temporal life among gods and men as a perpetual
struggle, which matters only that it be waged as nobly,
as vigorously and as wisely as possible, since all the real
peace that can fall to the divine people and the noble stock in
the course of time springs like a goddess of victory from the struggle;
this view must to a certain degree be found among
every people that is to assert its independence and
be allowed to develop its powers, even in the most peaceful
manner, and that it has not died out in the North either,
the present century has shown; but its spiritual stamp
it has almost lost, so that, instead of developing into
clarity, it was more and more obscured, until it became
the general opinion that wars are indeed a necessary
evil, because there are rapacious peoples, ambitious champions
and domineering princes; but that all struggle is nevertheless
at bottom abominable, and the more abominable, the more spiritual
its object is\dots{}. One may well think that it is
not my intention thereby to deify Asa-Thor or thus
to wash the blood from the champions of the North, so that they seem
angel-pure, for all who are versed in history can see that, just as
there has been only one hero with an unstained shield,
so also the champion spirit of Christianity is the
only one that deserves the name of divinity and worship; but
I should well like to impress upon all my readers the
misjudged truth that, although the earthly
instruments of the champion spirit were as far from being blameless as
the peacemakers, they were nevertheless neither less
indispensable nor less entitled to admiration and
indulgent judgment in a sinful world.''

With these words the decisive point of view is, insofar
% --- p. 424 ---
\setcounter{footnote}{0}%
\opage{424}as we can judge, stated once for all: mythology,
history and religion illuminate one another mutually.
--- One may disagree about the manner in which the
life-thought expressed can best be developed and put into
force, but about the thought itself all who want a little
more than a curious peeping into the kaleidoscope of the
combinations of the day must undoubtedly agree. The difference between
the school of knowledge and the popular school is here as striking
as possible. For when the emphasis is laid one-sidedly on
knowledge, mythology, history and religion must naturally
fall apart from one another into three subjects, three separate pieces
that have little or nothing to do with one another.

From mythology as a subject of knowledge one then learns a
good deal about Norse, Greek, Roman gods with accompanying
goddesses: it is many names, many figures,
attributes, scenes and fabulous exploits that the pupil has
to impress upon himself, if he is to find his bearings in the whole of it.
And what does all this matter to us; we do not
believe any longer in the old gods, do we? ``Certainly'' --- one
answers --- ``as far as the life of freedom, our inner life of the soul
with the world-view now prevailing, is concerned, the
whole thing really does not matter to us at all: we are far
beyond those old fantasies.'' And why must one
then apply oneself to knowing them? ``It is necessary
for the sake of cultivation; it is quite necessary to know
something of the ideal conceptions of antiquity. The mythic
poetic creations, in which the sublime and the terrible often
combine with traits of childlike grace and naive amiability,
contain, moreover, a rich variety of
artistic and poetic motifs. How should one
be able to understand the many allusions in the poets of the present day,
% --- p. 425 ---
\setcounter{footnote}{0}%
\opage{425}if one were not reasonably well oriented in one's mythology.
Anyone who does not know that Frigga was Odin's wife,
that Venus was married to Vulcan, that Mars was the god of war
among the Romans just as Thor was among the Norsemen, must surely
sit like an idiot in the company of cultivated people, cannot
visit a museum, cannot even go to the exhibition
and open his mouth without prostituting himself.''
Thus it is not the essence of myth, but only its semblance,
an aesthetically reflected afterglow, for whose sake the exertions
are made.

And now history? One has, as is well known, compendia
of the history of art, which give a survey of the
most important schools in the various periods of art, with
names of famous artists in each period together with a
short account of each artist's life, works,
birthplace, year of death, and so on. Now if a young person,
if not exactly from any special love of art, then
at least because it simply belongs to higher cultivation,
wished to acquaint himself somewhat with the history of art
and for that reason made a journey to one of the cities that
are distinguished by rich art treasures and famous galleries:
then this might very well be listened to. But if
he, after having arrived at a principal seat of art,
instead of walking about in the galleries and
contemplating the masterpieces, remained sitting in his lodgings
reading over the compendium until he knew it by heart,
and then, just as he was close to knowing it and thereby
at last prepared to enjoy the sight, discovered that
his time had run out, that he now, because of other
more important business, had to leave the city, then people would
surely find that he had taken the matter somewhat backwards.

% --- p. 426 ---
\setcounter{footnote}{0}%
\opage{426}What is self-evident in itself with regard
to the history of art does not yet seem to be quite
self-evident with regard to world history, for here it
is the compendium first and the compendium last on which instruction
turns. As a rule it costs so many
years of effort to learn more or less by heart an epitome of
world history well larded with names and dates,
that when the goal is reached, the time is up; then the
reverse movement sets in: one easily forgets
what one learned laboriously, and the hazy outline
that remains is still further from making the
impression of real history than the elephant in an ABC book
is from making the impression of a real elephant. The
conquest of a kingdom can be told in two lines, a whole revolution
in half a page; even if a leaf
or two is spent on describing a world upheaval, does
the reader thereby really enter into the real
events? One holds to the facts
in order to grasp the real, and does not notice that the reality
itself, with its strife and pain, its forces
and passions, has evaporated. What then do we learn from
history? ``Many things, remarkable things, an extraordinary
amount!'' And yet, when the question concerns the
present moment, an extraordinary little. As little as
the dead can write laws for the living, just as little
can the past immediately determine what
is to happen in the present. It is indeed true that
posterity always stands on the shoulders of the past; but it
is for that reason no less true that every age
with spirit and life has tasks with whose solution it must
begin afresh. ``What must be done?'' ``What will
% --- p. 427 ---
\setcounter{footnote}{0}%
\opage{427}happen?'' When these questions become pressing,
it avails little to look into a world-historical
compendium and ask the past for counsel; for the past's
oracles are obscure. It is with the past, when one
in an external, finite way has recourse to it for advice, as
with the witch of Endor: the utmost it can do
is to conjure up a ghost. Every bold undertaking
always has in the events of the past just as many grounds
dissuading as persuading; if the courageous man wants to dismiss
the deterrent examples by saying: ``that was
then, but now it is different'', then the cowardly
doubter can with the very same words invalidate any
encouraging example whatever. By holding one-sidedly
to the so-called pure facts, the naked data,
--- and that is what the compendium does on behalf of the items of knowledge ---
one grasps an externality without deeper understanding of the
corresponding internality. In such an externality one believes
one holds reality fast; that is a delusion.
The external real in the facts of history is
not the real itself, but only a picture of it, and not
even a complete picture, only some pale lines.
What shows itself is not; but it has been:
fuimus Troes, fuit Ilium et ingens gloria Teucrum. It
is the shadow of what has passed away that is taken for reality,
and then one imagines that this seeming reality is
the historical truth, the real history.

Is the separation between mythology and history
then to be based simply on this, that history gives us the real,
myth the unreal, and is nothing to count as
reality except what can be directly confirmed by
the testimony of the five senses: then it will be
% --- p. 428 ---
\setcounter{footnote}{0}%
\opage{428}difficult, not to say impossible, to form a certain
conviction about the facts of revealed religion.
For it cannot be hidden from an unbiased
eye that the biblical narratives contain two unlike
components, a supernatural and a natural one, of which
the first has a striking external likeness to the
mythical, the last to the straightforwardly historical; and
yet the phenomena of revelation lay claim to true
reality. According to the whole course of the development of knowledge, especially when we consider the ever-growing
influence of natural science, it is then also not strange that the chasm between
those who accept and those who reject positive
religion must grow ever greater. Every miracle
of which the Holy Scripture gives a description
naturally awakens in the reader the question:
is this real or is it unreal?

Those who, whether from lack of cultivation or
from lack of reflection or from other, perhaps deeper
personal, grounds, do not trouble themselves at all about nature's
so-called ``immutable laws'', but without any
intermediate determinations refer everything to God's almighty will,
can indeed without self-contradiction regard miracles
as real; but as concerns the opposition of the
natural and the supernatural, it will be impossible for them, so long
as they hold merely to the externals of the facts,
to specify any definite boundary
between the historical and the mythical. One reasons in general
thus: ``That Balaam's ass really was
able to speak is historical, for it stands in the Bible, but
that Achilles' horse Xanthos once is supposed to have had ``the
winged word'' is mythical, for it is not in the
% --- p. 429 ---
\setcounter{footnote}{0}%
\opage{429}Bible, it is only in Homer''. Such a piece of reasoning
can, however, hold only so long as it stands
firm that everything that stands in the Bible is externally
literally historical. And yet this presupposition,
on which Catholics and Protestants are officially agreed,
will not even suffice. The Catholics will, as is well known,
not admit that the age of miracles is past; they hold,
on the contrary, that in the history of the Church miracles have incessantly occurred
and still occur. Whoever, for example, cannot believe
that the pious Canon A ``walked dry-shod across
the Elbe'' or that the holy Bishop B was able to ``hang
his gloves on a sunbeam'', cannot believe either that Elijah went up to heaven in a glowing chariot, and that the walls
of Jericho fell because Joshua had the priests blow trumpets;
but whoever, despite ``the unanimous testimony of the holy Church'',
cannot believe such things is ``a heathen'',
``an unbeliever'', ``a heretic'', ``a Turk''!

Those who hold exclusively to the likeness, mark you,
the external likeness, between the mythical in heathenism
and the supernatural in Christianity, can,
if they want to be consistent, not rest until
all the supernatural has become mythical, and the whole
positive religion been transformed into mythology.

There is one important point on which atheists and theologians
are always agreed, namely the externally historical.
``If what the Bible reports is true, then it must have really
happened; if it has really happened, then it must be
historical'': that is the common presupposition; on
the externally historical reality, truth is to rest.
When theology nevertheless, as ``historical science'',
% --- p. 430 ---
\setcounter{footnote}{0}%
\opage{430}will, as concerns holy history, come to a
different conclusion than atheism, it cannot shake
the major premise, but must cautiously confine itself to
tinkering with the minor premise.

``If the supernatural is really true, then it must
be historical.

Now it is not historical and cannot be so.

Ergo the supernatural is in reality untrue,
that is, it is mythical''.

Thus the conclusion is atheistic.

``If the supernatural is really true, then
it must be historical.

Now it is in all probability really historical.

Ergo it is in every way certain and true, and the
mythical excluded.''

Thus the conclusion is theological.

It is the externally historical, the scriptural-historical,
over which apologetics and negative criticism are incessantly
pulling at one another. And in that contest apologetics always draws
the shortest straw; understand, it is scientific
straw, otherwise the apologetic game could not be continued.
The first of the cited conclusions is at least,
as regards form, logically correct; the
last, on the other hand, is lame, incorrect\footnote[1]{After men such as G. E. Lessing and S. Kierkegaard have sufficiently demonstrated the impossibility of ``building one's salvation'' on objective certainty about anything objectively historical, apologetics might well spare itself the exertion of benefiting religion by patching together ``synopses'' and mending ``holy history''.}. One needs
only to bethink oneself of the difference between historical and
% --- p. 431 ---
\setcounter{footnote}{0}%
\opage{431}religious faith, between probability and certainty, and one
will get an idea of how matters stand
when theology, in order to save a few historical scraps,
must give way to criticism and, in order to preserve a semblance of science,
in the name of knowledge stoop to flirting
with a slightly veiled atheism.

According to the point of view stated by Grundtvig, the relation
between truth and reality will show itself in a
different light. Reality is two-sided; it has not
merely an externality, but also an internality. It is
so far from the case that the emphasis can fall on the externality,
that, where the question is about truth,
it falls much rather on the internality. This holds
just as decisively for mythology as for history; for
the internal real is spirit. The same spirit
that is in a people's mythology must also be in its history,
and conversely. History and mythology must reciprocally
clarify and illuminate each other, if it is otherwise to
be possible to determine the spirit of the people. The history of antiquity
is at once a history of gods and of heroes:
such heroes, such gods, and then again:
such gods, such heroes!

Just as the heroic life, the higher human life, has its prototype
in the life of the gods, and only under this presupposition
can be real, so in turn the life of the gods
has its natural ground and the force of its reality in the heroic life.
The difference between Greek and Nordic
mythology is in reality itself a difference between
the Greek and the Nordic spirit of the people, but the difference
between the Greek and the Nordic spirit of the people is
% --- p. 432 ---
\setcounter{footnote}{0}%
\opage{432}precisely the fundamental cause from which alone what is characteristic in
the history of both peoples can be derived\footnote[1]{The author of this essay has sought to show this in the lectures: On Obstacles and Conditions (second lecture).}.

It is with a whole people as with a single
human being: in order to gather itself for heightened self-development,
it must know itself, its forces and natural dispositions.
The testimony of its natural dispositions and of the direction of life
conditioned by them is to be sought in its mythology; the use
it has hitherto made of forces and dispositions must be seen
from its course of life, its history. The best means of
getting away from the confusion of the moment and bethinking oneself of
the future is to go back to the past, but, mark you,
along the straight road of spirit, not along the manifold winding, often
bewildering paths of external finite events.
By losing oneself in the external manifold the mind is distracted,
the gaze goes astray, and the contemplation ends in
stupefaction; only when the manifoldness of events is governed
by the unity of thought, so that every development
both proceeds from and leads to the answering of a practical
main question, will history contribute to the consciousness
becoming clear, to the people really going into itself.
Here in the past there is then much that strengthens and uplifts,
but also much that humbles and depresses, errors that
must be corrected, blunders and neglects that must be repented,
weaknesses that must be combated, blows of fate that
must be borne with resignation. But in order that history in
this way may be able to awaken and instruct,
the external reality must everywhere be referred to and explained
by the inner. History is then not something merely
lived out, a burial place for vanished glory, the
% --- p. 433 ---
\setcounter{footnote}{0}%
\opage{433}past's dead sea, where memory alone dwells among
sunken cities and thrones, but a present,
moving power, a stream of life, the river of the ages
that drives the wheel of existence.

Only when the spiritual bond that unites mythology
and history is found can religion be cognized in
its peculiar relation to both. The fundamental concept
to which the reciprocal action of mythology, history and religion
must be referred is the concept of freedom. The principle of freedom
is will, the essence of will is spirit; in the world of freedom
it is spirit and will that govern and determine
the external reality. If the spirit of the people is awake and the will
strong, then the power to act is great, and history rich
in deeds; if the spirit is drowsy and the will feeble, then
the power to act is slight, and the real history insignificant.
Now since mythology is the mirror of the spirit in which
the people sees its life-ideals and aims, so
this first natural self-revelation of the spirit of the people
at the same time indicates the height to which a people by its own forces
is able to raise itself, the measure of freedom to which
its real history can attain, and in which it naturally
must have its boundary. Thus the Greeks
in the ancient world attained precisely to the realization of the
ideal of freedom which in the form of beauty expresses itself
in their mythology; but not a step further.
What holds of the ideal of beauty among the Greeks
holds in a corresponding way of the ideal of power among the Romans.
Seen externally, it is the arms of the Barbarians that
put an end to the power of Rome; but, seen internally, this
power had already from its beginning its goal, its
boundary in itself.

% --- p. 434 ---
\setcounter{footnote}{0}%
\opage{434}If now in reality there were no other
spirit than the spirit of the people, then it would look sad
for freedom. For every spirit of a people has namely in the people's
individuality a natural boundary that cannot be transgressed.
To be sure spirit as spirit is always the universal, but
this universal gets reality, with flesh and blood,
only by living and dwelling in the individual. To want to free
spirit from the natural limitation by dissolving
the national peculiarity in cosmopolitan uniformity
is, as experience has amply shown, a metaphysical
error, an attempt to put the shadow in the
place of the body and to make life itself a bloodless abstraction.

It lies in the nature of the case that an increasing civilization
and with it advancing thought must gradually
undermine the mythical popular faith and transform
the figures of the gods into conscious symbols. If now a
mythology could be thought to contain the highest truth
and to state the principle of real freedom: then
the transformation that the myths underwent in reflection
would surely signify not a weakening, but an elevation
of the spirit of the people, which in this way came to
clarity about itself. But this is not how matters
stand. It is not only the external reality of the gods
that the critique of the understanding attacks; it attacks the essence
of the world of the gods, its truth, as reason sees that such an
essence no longer promotes, but hinders freedom.
Then the dissolution follows as of itself; for the
conscious thinking that discovers the irrational, and with it
the spiritless, in the doctrine of the gods, is not able to install
a new spirit in place of the old. Thus could
% --- p. 435 ---
\setcounter{footnote}{0}%
\opage{435}the Greek philosophy well dissolve the Greek mythology,
but the dissolution attacked the spirit of the people; philosophy itself
was magically governed by the basic forms of the spirit of the people:
by dissolving mythology it reached its perfection,
but just as it was on the point of going further,
it initiated its own dissolution, for the boundary was set
by spirit.

With revealed religion comes the spirit of truth,
the spirit of will and of freedom, which alone is able
to solve the riddles of reality. The relation of religion
to mythology is from this standpoint to be determined
by the relation of the revealed spirit to the spirit of the people.
To see what lies in this, we can begin
with the single human being. Revealed religion
is raised above all national limitation; it is destined
for all peoples, hence for all human beings, for
every single human being without exception. But although
every human being must be able to appropriate religion, there
is nevertheless, as concerns the depth, inwardness, force and independence
of such an appropriation, an unmistakable
difference between individuals. Individuality
is more than a passive medium for the passage of spirit,
it is like the gem that refracts the light and
throws back the rays with a peculiar luster and
play of colors. It can be said without exaggeration that every
true and capable individuality grasps the religious truths
in a peculiar way. The difference does not derive
directly from the externals, for the confession can
be the same and the doctrine the same; the essential
difference is in the internal. Not only in the external,
but also in the internal a man is
% --- p. 436 ---
\setcounter{footnote}{0}%
\opage{436}different from a woman; from which in turn it follows that the religious
life assumes another shape in the man than in the woman.
If revealed religion did not lay the
very greatest weight on the preservation and perpetuation of what is individual
in the life of freedom, the doctrine of the resurrection of the dead would
be without significance. --- But the individuality in
the single individual is on the whole inseparable from the individuality of the people.
That in the latter as well as in the former there is much
that belongs to finitude and falls prey to perishability
shall not be denied --- this holds above all
of the external --- but in the name of the general, individuality
must be upheld. If now the significance of the individuality of the people
is held fast, then there is also nothing unreasonable
in one people's being disposed to a deeper conception
of the life of freedom than another; Nordic mythology
in this respect bears witness to the peculiar
endowment of the Nordic peoples. To assert a popular
advantage is by no means the same as
flattering the vanity of the people; on the contrary it can be
an attempt to awaken its conscience. Here, as everywhere, it holds:
``From him to whom much is entrusted,
much may also be required''. When a great past stands
in tragic conflict with a wretched present, then something
is in the way, something has been lost. History bears witness
against the people in question, that it has neglected the
opportune time and not used its great forces.

The deeper a mythology is, the deeper is the spirit of the people,
and the deeper the spirit of the people is, the more inwardly
and forcefully will the people be able to appropriate the spirit
of revelation. It is the life of freedom, hence reality,
the inner reality, on which everything first and last
% --- p. 437 ---
\setcounter{footnote}{0}%
\opage{437}depends. Whoever cannot understand that the mythical
symbols to this very day have something to do with real
life will, when we look more closely, also
be unable to feel and grasp that the supernatural
facts of revelation have something to do with the spiritual
conditions of the present: the truth, the inner reality
in the former as in the latter, is spiritual reality, reality
for spirit. But just as there is a difference between temporal
and eternal life, between temporal and eternal freedom, between man's
self-will and God's almighty, holy will, so
there is also a difference between the spirits of the peoples and the holy spirit of
religion; and hence also a difference as wide as heaven
between the idea-reality that can be ascribed to
the mythical symbols and the ideal reality that
belongs to the supernatural, a difference in principle between
mythology and revelation\footnote[1]{On religion and mythology cf. the author's Philosophy of Religion p. 11 ff.}.

The fundamental thought which is here indicated in fleeting outline
should now in the popular school assume shape, life and
color; for this teaching forces are needed. It already costs
exertion of spirit and sustained thinking to work
one's way many-sidedly and thoroughly into a rich world-view, and
yet this is still only half the way to the goal; for the
difficulty is to get what has been gained unfolded and
presented. We must on this occasion call attention
to the categorical difference there is between
lecturing, chatting genially, and speaking personally.

To lecture is to present a subject objectively, to clarify
its connection in general concepts and to show its validity
by virtue of incontrovertible grounds. In the strict

% --- p. 438 ---
\setcounter{footnote}{0}%
\opage{438}academic lecture the teacher holds back his subjectivity;
if in individual utterances he lets it shine through,
it happens so indirectly that the passing gleams of
individual mood serve only to strengthen the impression
of the significance of the subject and to sharpen attention to
the weight of the grounds. Only by putting a check on his
individuality and removing irrelevant regard for private
opinions can the teacher bring about the exalted
calm, the cool composure, the undimmed clarity
of consciousness, that is a condition for the listener
to be able undisturbed to follow the movement of thought point by
point.

To chat genially is now precisely the opposite. The
natural, fresh, immediate subjectivity does not creep into
hiding behind the objective, but gives itself quite as
it is, unconcerned whether it happens in Sunday dress or
\textit{in puris naturalibus} or in dressing gown and slippers. The
genial discourse can play in all colors, strum on
all strings and produce a highly varied effect. It
can be attractive and entertaining, if the speaking
or, to stay in the category, chatting individuality
has whim, wit and humor in abundance; but if these
good gifts are lacking, it can also become long-drawn-out,
tiresome and boring. It can become something instructive,
if he who playfully plays with the word has
insight, vision and understanding, but in the opposite case it
can also become something else, for example nonsense. Over
all this a higher, unfathomable contingency rules.

He who wants to speak personally must now neither keep
to the coldness of objectivity nor warm himself at the blaze of
the fireworks of genial whims. He may not, as regards the
% --- p. 439 ---
\setcounter{footnote}{0}%
\opage{439}objective, offer his listeners instruction coming from outside
and let them try to put on
a theory that is not particularly for them but for
humanity, just as one lets the recruit try on a
uniform that is not sewn for him but for the army;
first of all, if we may say so, he must have taken the measure
of himself and his listeners, so that he is really able
to put himself at the listeners' standpoint, to imagine
himself into their circle of ideas, to speak out of their own
consciousness, to touch the strings in their own breast and to help
them think their own thoughts. This is the task,
seen from the subjective side. But the subject matter, the objective,
is just as important to him. In this respect he must
bring along a fullness of knowledge, a rich
material, and have the material duly worked through. It
is not the theme taken by itself, any more than it is the form of
merely individual peculiarity, that gives the speech
the personal stamp: it is the spiritual unity of
both. The teacher must have taken the subject so far
into his consciousness that he literally lives in it. The
whole must be so present to him that in the moment
of speaking, in the course of the lecture, according to thought and
mood he can freely seize on the examples, can master and
shape the manifold in such a way that a popular character-portrait,
a painting of souls from the real life of history,
involuntarily forms itself for the listeners. Then the
highest becomes simple and is understood simply; thought is clothed in
narrative; images and reflections, examples and
descriptions alternate with full freedom, but without circumlocution
or digressive prolixity: thus the whole
presentation works to secure understanding and promote
% --- p. 440 ---
\setcounter{footnote}{0}%
\opage{440}cognition's development, then the unlearned can naturally
follow along, and the learned take pleasure in sitting
by and listening. Mere liveliness cannot do
this, any more than dry thoroughness. In the domain of the doctrine of spirit
one is all too inclined to set a divorce
between the thorough and the awakening lecture; in that one assumes that the thorough belongs exclusively
to science and therefore cannot condescend
to be awakening, whereas the awakening, which falls
outside science, cannot possibly be thorough, but
must be content to be lively. That is a false point of view.
The truth is that a popular lecture must be
given in another form than a scientific one; but those
who assume that a lecture is unthorough because it is
awakening and living must presumably also assume that a
lecture is always scientific and thorough when it
is merely properly sluggish, untalented and tedious.

From what has been treated here it will, it is hoped, now be
clear that the progress made in the new popular schools
does not rest on the prior knowledge of the pupils, but
on the competence of the teachers. In the school of knowledge the
predominant part of the burden of work falls on the pupils, for it
is really easier to hear a lesson than to learn
a lesson; in the popular school, on the other hand, the burden of work
falls chiefly on the teacher, inasmuch as it is always easier
to listen to a fruitfully awakening lecture than it is
to give it. If only the popular lecture can develop
to the clarity and perfection that the importance of the subject
demands, then one may be assured that the listeners
will surely get to work, that every possible
guidance, be it oral or written, be it
% --- p. 441 ---
\setcounter{footnote}{0}%
\opage{441}to be found in smaller textbooks or in larger works, will
be used as time and occasion allow, for enterprising self-activity
with fresh forces accomplishes great things.

It has not been our task in the present essay
either to revise given plans of instruction or
to draw up outlines for new ones; we have, with special reference to
idea and reality, confined ourselves to directing attention
to a great undertaking, which, thanks to
the many young men who not only among us but also
in the sister countries have with self-sacrificing enthusiasm laid
hand to the work, will with the assistance of the government bear
fruit for spiritual freedom and civic progress,
to the good fortune and blessing of the lands of the North. We have
fixed our gaze on the awakening on which, in our conviction,
alone a hope for the future can be built;
but we have not wished either to overlook or to misjudge
the good and useful that the school of knowledge,
in lower as well as higher style, has rendered and will continue
to render: we have, with the difference in principle between
learned and popular cultivation in view, wished to contribute to a
cognition that in the present there can and must be
two kinds of instruction.
\signature{R. \emph{Nielsen}.}

\end{document}
